% Part: second-order-logic
% Chapter: syntax-and-semantics
% Section: expressive-power

\documentclass[../../../include/open-logic-section]{subfiles}

\begin{document}

\olfileid[hi]{sol}{syn}{exp}
\olsection{अभिव्यंजक शक्ति}

\begin{explain}
द्वितीय-कोटि चरों पर परिमाणीकरण भाषा की अभिव्यंजक शक्ति को प्रथम-कोटि
तर्कशास्त्र की अपेक्षा अत्यधिक बढ़ा देता है। द्वितीय-कोटि अस्तित्ववाचक
परिमाणीकरण हमें यह कहने देता है कि कुछ गुणों वाले फलन या संबंध अस्तित्व में
हैं। प्रथम-कोटि तर्कशास्त्र में ऐसा करने का एकमात्र उपाय इस प्रयोजन के लिए
कोई अतार्किक चिह्न (अर्थात् !!a{function} या !!{predicate}) विनिर्दिष्ट करना
है। द्वितीय-कोटि सार्विक परिमाणीकरण हमें यह कहने देता है कि !!{domain} के
सभी उपसमुच्चयों, उस पर सभी संबंधों अथवा !!{domain} से !!{domain} में जाने वाले
सभी फलनों में कोई गुण है। प्रथम-कोटि तर्कशास्त्र में हम केवल यह कह सकते हैं
कि भाषा के किसी अतार्किक चिह्न को दिए गए उपसमुच्चय, संबंध या फलन में वह गुण
है। और जब हम कहते हैं कि किसी गुण वाले उपसमुच्चय, संबंध या फलन अस्तित्व में
हैं, अथवा उन सभी में वह गुण है, तो उस गुण को विनिर्दिष्ट करने में भी
द्वितीय-कोटि परिमाणीकरण का उपयोग कर सकते हैं। इससे हम ऐसे संबंधों को परिभाषित
कर सकते हैं जो प्रथम-कोटि तर्कशास्त्र में परिभाषणीय नहीं हैं, और प्रांत के
ऐसे गुण व्यक्त कर सकते हैं जो प्रथम-कोटि तर्कशास्त्र में अभिव्यक्त नहीं होते।
\end{explain}

\begin{defn}
यदि $\Struct{M}$ भाषा~$\Lang{L}$ की !!a{structure} है, तो संबंध
~$R \subseteq \Domain{M}^2$, भाषा~$\Lang{L}$ में \emph{परिभाषणीय} है यदि
कोई ऐसा !!{formula}~$!A_R(x, y)$ हो जिसमें केवल $x$ और~$y$ मुक्त हों और
$R(a, b)$ लागू हो (अर्थात् $\tuple{a,b} \in R$) तभी और केवल तभी जब
$s(x) = a$ और $s(y) = b$ के लिए $\Sat{M}{!A_R(x, y)}[s]$।
\end{defn}

\begin{ex}
प्रथम-कोटि तर्कशास्त्र में हम तादात्म्य संबंध~$\Id{\Domain{M}}$ (अर्थात्
$\Setabs{\tuple{a,a}}{a \in \Domain{M}}$) को सूत्र $\eq[x][y]$ से
परिभाषित कर सकते हैं। द्वितीय-कोटि तर्कशास्त्र में हम इस संबंध को
\emph{बिना~$\eq$ के} परिभाषित कर सकते हैं। यदि $a$ और $b$,
~$\Domain{M}$ के वही !!{element} हैं, तो वे~$\Domain{M}$ के उन्हीं
उपसमुच्चयों के !!{element} हैं (क्योंकि समुच्चय अपने !!{element} से निर्धारित
होते हैं)। इसके विपरीत, यदि $a$ और $b$ भिन्न हों, तो वे उन्हीं उपसमुच्चयों
के !!{element} नहीं हैं: उदाहरणार्थ, यदि $a \neq b$, तो $a \in \{a\}$ किंतु
$b \notin \{a\}$। अतः ``~$\Domain{M}$ के उन्हीं उपसमुच्चयों का !!{element}
होना'' ऐसा संबंध है जो $a$ और $b$ पर तभी और केवल तभी लागू होता है जब
$a = b$। यह द्वितीय-कोटि तर्कशास्त्र में व्यक्त होने वाला संबंध है, क्योंकि
हम~$\Domain{M}$ के सभी उपसमुच्चयों पर परिमाणीकरण कर सकते हैं। अतः निम्न
!!{formula}, $\Id{\Domain{M}}$ को परिभाषित करता है:
\[
\lforall[X][(X(x) \liff X(y))]
\]
\end{ex}

\begin{prob}
दिखाएँ कि $\lforall[X][(X(x) \lif X(y))]$ (ध्यान दें: $\liff$ नहीं,
बल्कि~$\lif$!) $\Id{\Domain{M}}$ को परिभाषित करता है।
\end{prob}

\begin{ex}
यदि $R$ द्विस्थानी !!{predicate} है, तो $\Assign{R}{M}$,
~$\Domain{M}$ पर द्विस्थानी संबंध है। कुछ भ्रम की संभावना के बावजूद, हम
$R$ का प्रयोग~$R$ के !!{predicate} और स्वयं संबंध~$\Assign{R}{M}$, दोनों के
लिए करेंगे। $R$ का \emph{संक्रामक संवरण}~$R^*$ वह संबंध है जो $a$ और $b$
के बीच तभी और केवल तभी लागू होता है जब किन्हीं $c_1$, \dots, $c_k$ के लिए
$R(a,c_1)$, $R(c_1, c_2)$, \dots, $R(c_k,b)$ लागू हों। इसमें $k = 0$ की
स्थिति भी आती है, अर्थात् यदि $R(a,b)$ लागू है तो $R^*(a,b)$ भी लागू है।
इसका अर्थ है कि $R \subseteq R^*$। वास्तव में $R^*$ वह सबसे छोटा संबंध है
जिसमें~$R$ सम्मिलित है और जो संक्रामक है। द्वितीय-कोटि तर्कशास्त्र में हम
कह सकते हैं कि $X$ ऐसा संक्रामक संबंध है जिसमें~$R$ सम्मिलित है:
\begin{multline*}
  !B_R(X) \ident \lforall[x][\lforall[y][(R(x,y) \lif X(x, y))]] \land {}\\
\lforall[x][\lforall[y][\lforall[z][((X(x,y) \land X(y,z)) \lif X(x,
      z))]]].
\end{multline*}
पहला संयोज्य कहता है कि $R \subseteq X$, और दूसरा कि $X$ संक्रामक है।

यह कहना कि $X$ ऐसा सबसे छोटा संबंध है, यह कहना है कि $X$ स्वयं प्रत्येक उस
संबंध में सम्मिलित है जिसमें $R$ सम्मिलित है और जो संक्रामक है। अतः हम
~$R$ के संक्रामक संवरण को निम्न !!{formula} से परिभाषित कर सकते हैं:
\[
R^*(X) \ident !B_R(X) \land \lforall[Y][(!B_R(Y) \lif
  \lforall[x][\lforall[y][(X(x, y) \lif Y(x,y))]])].
\]
हमारे पास $\Sat{M}{R^*(X)}[s]$ तभी और केवल तभी जब $s(X) = R^*$।
~$R$ का संक्रामक संवरण प्रथम-कोटि तर्कशास्त्र में व्यक्त नहीं किया जा सकता।
\end{ex}

\end{document}
